% !TeX root = ../main.tex
\chapter{Functional Specifications}
\label{app:functional_specs}
The functional specifications used for the design are summarized in
Table~\ref{tab:functional-specifications}. These values define the intended
operating envelope of the prototype and establish the quantities that must be
checked during circuit analysis and later verification testing.

\begin{table}[H]
    \centering
    \caption{Functional specifications.}
    \label{tab:functional-specifications}
    \begin{tabularx}{\textwidth}{>{\raggedright\arraybackslash}p{0.27\textwidth}X}
        \toprule
        \textbf{Function} & \textbf{Specification} \\
        \midrule
        Monitored rails & +5 V A, +5 V B, +5 V, -5 V, +15 V, and -15 V \\
        Number of channels & Six independent voltage/current measurement channels \\
        +5 V A/B current range & 0 to 2 A design range per rail \\
        +5 V, -5 V, +15 V, -15 V current range & 0 to 1 A design range per rail \\
        Voltage measurement resolution & 25 mV/LSB nominal using the LTC2945 rail-voltage ADC \\
        Current-sense resolution & 25 \(\mu\)V/LSB across the selected shunt resistor; approximately 0.641 mA/LSB with 39 m\(\Omega\) and 0.500 mA/LSB with 50 m\(\Omega\) \\
        Isolation & Independent digital isolation for every floating monitor channel \\
        Alarm processing & User-configurable upper/lower voltage and current limits with local and EPICS status indication \\
        Fault handling & Detect failed I\textsuperscript{2}C transactions and prevent invalid data from being treated as a valid measurement \\
        Local interface & GIGA Display Shield showing voltage, current, alarm state, communication state, and configured limits \\
        Network interface & Ethernet PSC server for EPICS IOC communication plus browser-based configuration/monitoring services \\
        Configuration retention & Network parameters and alarm limits stored in nonvolatile memory \\
        Controller logic level & 3.3 V; external interfaces must not exceed the GIGA I/O limits \\
        \bottomrule
    \end{tabularx}
\end{table}